\documentclass[12pt,a4paper]{article}
\usepackage{amsmath,amssymb,amsthm,mathtools}
\usepackage{mathrsfs}
\usepackage{setspace,graphicx,array,tikz,url}
\usepackage{cite}
\usepackage[hidelinks]{hyperref}
\usepackage{fancyhdr}
\usepackage{lastpage}
\usepackage[top=3cm,left=3cm,right=3cm,bottom=3cm]{geometry}
\usepackage{enumitem}
\usepackage{titlesec}
\usepackage{etoolbox}
\usepackage{microtype}

\patchcmd{\thebibliography}
  {\advance\leftmargin\labelsep}
  {\labelsep=.5em\leftmargin=2.5em\itemindent=0em\advance\leftmargin\labelsep}
  {}{}

\titlespacing*{\section}{0pt}{3em}{2em}

\setstretch{1.1}

\pagestyle{fancy}
\fancyhf{} % Clear header and footer
\fancyfoot[C]{\footnotesize \thepage\ of \pageref{LastPage}}

\renewcommand{\headrulewidth}{0pt}
\renewcommand{\qedsymbol}{$\blacksquare$}

\newtheoremstyle{spaced}
  {15pt}   % Space above (e.g., 10pt)
  {20pt}   % Space below (e.g., 10pt)
  {\itshape} % Body font (e.g., italic)
  {}       % Indentation of the header
  {\bfseries} % Theorem header font (e.g., bold)
  {.}      % Punctuation after theorem header
  {.5em}   % Space after theorem header
  {}       % Theorem head spec (can be left empty)

\theoremstyle{spaced}
\newtheorem{conjecture}{Conjecture}[section]
\newtheorem{theorem}{Theorem}[section]
\newtheorem{definition}{Definition}[section]
\newtheorem{example}{Example}[section]
\newtheorem{lemma}{Lemma}[section]
\newtheorem{proposition}{Proposition}[section]
\newtheorem{corollary}{Corollary}[section]
\newtheorem{remark}{Remark}[section]

\title{Banach Tarski Paradox}
\author{
  Mohammad Javad Moghaddas Mehr\thanks{m.moghadas11235@gmail.com}}

\date{\today}

\begin{document}
\maketitle
\vspace{5em}
\begin{abstract}
	This paper gives a self-contained proof of the Banach--Tarski paradox. We outline the foundational context, develop the required group-theoretic machinery, construct a paradoxical decomposition of the sphere, and extend it to a solid ball using equidecomposability. The proof illustrates how free group actions and the Axiom of Choice combine to produce the classical paradox.
\end{abstract}

\small{\textbf{Keywords:} Banach–Tarski Paradox, Axiom of Choice, Measure Theory, Paradoxical Decomposition}


\section{Introduction}

The problem of assigning a meaningful notion of ``size'' to arbitrary subsets of
$\mathbb{R}^n$ lies at the heart of measure theory and modern analysis. Classical
intuition suggests that such a notion should extend the familiar geometric
concepts of length, area, and volume, and should respect the underlying
symmetries of Euclidean space. One is therefore naturally led to seek a measure
$\mu$ defined on all subsets of $\mathbb{R}^n$ that satisfies the following basic
requirements:
\begin{enumerate}
	\item \emph{countable additivity} on disjoint families;
	\item \emph{invariance} under rigid motions;
	\item \emph{normalization} on standard sets, such as $\mu([0,1]^n)=1$.
\end{enumerate}

A classical argument shows that these requirements are incompatible.
By using the axiom of choice to select one representative from each orbit under
rational translation, one can partition the unit interval in a way that forces any
countably additive, translation-invariant measure to assign contradictory values.
Thus no such measure can exist on all subsets of $\mathbb{R}$.

In higher dimensions an even more striking phenomenon occurs. In 1924 Banach and
Tarski~\cite{banach1924decomposition} showed that in $\mathbb{R}^n$ for $n\ge 3$,
one may dispense with countable additivity entirely and still encounter
impossibility. They proved that any two bounded open sets in $\mathbb{R}^3$ are
\emph{equidecomposable}: each can be partitioned into finitely many pieces that
can be rearranged by rigid motions to form the other. In particular, a solid
ball can be divided into finitely many disjoint, highly nonmeasurable pieces,
which may then be reassembled into two balls congruent to the original. A sharp
result due to Robinson~\cite{robinson1947decomposition} shows that such a
doubling of the ball can be accomplished with \emph{exactly} five, and fewer than five pieces will not suffice.

At the conceptual core of the Banach--Tarski paradox is the idea of a
\emph{paradoxical decomposition}: a set that can be partitioned into finitely
many pieces which, after applying suitable isometries, form two disjoint copies
of the original set. Tarski later demonstrated that such decompositions occur
precisely for group actions containing a nonabelian free subgroup. This insight
places group theory—particularly the behavior of free groups and their
actions—at the center of the phenomenon. The classical proof constructs a
paradoxical decomposition of the free group $F_2$, embeds $F_2$ into
$\mathrm{SO}(3)$, and uses this embedding to transfer the paradox to the unit
sphere $S^2$. A radial extension then propagates the decomposition to the entire
unit ball.

The purpose of this report is to give a streamlined, self-contained exposition of
this strategy. In Section~2 we review the relevant concepts from group actions,
equidecomposability, and paradoxical decompositions. Section~3 develops the
proof: we obtain a paradoxical decomposition of the free group $F_2$, embed it
into $\mathrm{SO}(3)$, transfer the resulting decomposition to $S^2$, and extend
it to the solid unit ball via radial invariance. The presentation emphasizes the
interplay between group-theoretic structure and the axiom of choice, showing how
these ingredients combine to yield one of the most remarkable results in modern
set theory and geometric analysis.

\section{Preliminaries}
In this section we introduce the basic notions required for the Banach--Tarski paradox.
Throughout, $G$ denotes a group and $X$ a set.
Our goal is to understand how group actions allow us to ``move pieces around,'' and how
such motions lead to equidecompositions and paradoxical decompositions.

\begin{definition}[Group Action]
	A \emph{(left) action} of $G$ on $X$ is a function
	\[
		G \times X \longrightarrow X, \qquad (g,x) \mapsto g \cdot x,
	\]
	such that for all $g,h \in G$ and all $x \in X$,
	\begin{enumerate}
		\item $e \cdot x = x$, where $e$ is the identity element of $G$;
		\item $(gh)\cdot x = g \cdot (h \cdot x)$.
	\end{enumerate}
	In this situation we say that $G$ \emph{acts} on $X$.
\end{definition}

\begin{remark}
	A \emph{right} action is defined similarly using a function
	$X \times G \to X$ and the rule
	\((x \cdot g)\cdot h = x \cdot (gh)\).
\end{remark}

\begin{definition}[Equidecomposability]
	Let $G$ act on a set $X$.
	Two subsets $A,B \subseteq X$ are \emph{$G$-equidecomposable} if there exist
	pairwise disjoint subsets $A_1,\dots,A_n \subseteq A$ and elements
	$g_1,\dots,g_n \in G$ such that
	\[
		B = \bigcup_{i=1}^n g_i \cdot A_i.
	\]
	In this case we write $A \sim_G B$ and say that $A$ can be partitioned into finitely
	many pieces and moved by elements of $G$ to form $B$.
\end{definition}

\begin{definition}[Paradoxical Decomposition]
	Let $G$ act on $X$.
	A subset $A \subseteq X$ admits a \emph{paradoxical decomposition} if there exist
	pairwise disjoint subsets
	\[
		A_1,\dots,A_m, \qquad B_1,\dots,B_n \subseteq A
	\]
	and elements
	\[
		g_1,\dots,g_m, \qquad h_1,\dots,h_n \in G
	\]
	such that
	\[
		A = \bigcup_{i=1}^m g_i \cdot A_i
		\qquad\text{and}\qquad
		A = \bigcup_{j=1}^n h_j \cdot B_j.
	\]
	That is, $A$ can be partitioned into finitely many pieces which, after applying
	elements of $G$, form two disjoint copies of $A$.
\end{definition}

\begin{lemma}\label{lem:equidecomp_preserves_paradoxical}
	If $A \sim_G B$ and $A$ admits a paradoxical decomposition under $G$, then
	$B$ also admits a paradoxical decomposition under $G$.
\end{lemma}

\begin{proof}
	Suppose
	\[
		A = \bigcup_{i=1}^m g_i A_i = \bigcup_{j=1}^n h_j B_j
	\]
	is paradoxical, and that $A \sim_G B$ via a decomposition
	\[
		B = \bigcup_{k=1}^r t_k C_k, \qquad C_k \subseteq A.
	\]
	Replacing each $A_i$ and $B_j$ by its intersections with the $C_k$ (and then applying
	the corresponding $t_k$) yields a paradoxical decomposition of $B$.
\end{proof}

\begin{remark}\label{rem:countable_sets harmless}
	Removing a countable subset does not affect equidecomposability or paradoxicality.
	If $D \subset X$ is countable and $G$ acts on $X$, then $X$ is $G$-equidecomposable
	with $X \setminus D$.
	This follows from the fact that a countable set can be ``absorbed'' into a single
	piece of any finite equidecomposition by sending its points into an infinite orbit.
\end{remark}

A paradoxical decomposition is the formal mechanism behind the Banach--Tarski paradox.
In Theorem~\ref{thm:bt} we will show that the closed unit ball
\(
B = \{ x \in \mathbb{R}^3 : \|x\| \le 1 \}
\)
admits a paradoxical decomposition under the action of the group of rotations.
In particular, $B$ is equidecomposable with two disjoint copies of itself.

\section{The Banach--Tarski Paradox}

In this section we give a detailed proof of the Banach--Tarski paradox following the
classical strategy:
\begin{enumerate}
	\item construct a paradoxical decomposition of $F_2$; And realize this free group as a subgroup of the rotation group $\mathrm{SO}(3)$;
	\item transfer the paradoxical decomposition from $F_2$ to the unit sphere $S^2$;
	\item extend the paradoxical decomposition of $S^2$ to the solid unit ball in $\mathbb{R}^3$.
\end{enumerate}

Throughout, $G$-equidecomposability and paradoxical decomposition are meant in the
sense introduced in Section~2.

\subsection{A paradoxical decomposition of the free group $F_2$}
Let $F_2$ be the free group on generators $a$ and $b$, equipped with its standard
description in terms of reduced words.

\begin{definition}
	A word \(w = x_1 x_2 \cdots x_n\) in the alphabet
	\(\{a, a^{-1}, b, b^{-1}\}\) is \emph{reduced} if no letter is adjacent to its
	inverse.
	The free group \(F_2\) is the set of reduced words together with the
	empty word \(e\), with multiplication given by concatenation followed by
	cancellation of all occurrences of \(x x^{-1}\) and \(x^{-1} x\).
\end{definition}

Every nontrivial element of \(F_2\) begins with exactly one of the letters
\(a, a^{-1}, b,\) or \(b^{-1}\).  Define the corresponding subsets:
\[
	\begin{aligned}
		S(a)      & = \{\, w \ne e : w \text{ begins with } a \,\},      \\
		S(a^{-1}) & = \{\, w \ne e : w \text{ begins with } a^{-1} \,\}, \\
		S(b)      & = \{\, w \ne e : w \text{ begins with } b \,\},      \\
		S(b^{-1}) & = \{\, w \ne e : w \text{ begins with } b^{-1} \,\}.
	\end{aligned}
\]
These four sets are disjoint and partition all nontrivial words, giving
\begin{equation}\label{eq:F2partition-clean}
	F_2 = \{e\} \,\cup\, S(a) \,\cup\, S(a^{-1}) \,\cup\, S(b) \,\cup\, S(b^{-1}).
\end{equation}

For any \(w \in S(a^{-1})\) we may write \(w = a^{-1} u\), so
\(
aw = aa^{-1} u \to u
\)
after reduction.
Conversely, every reduced word not beginning with \(a\) arises uniquely in this way.
Applying the same reasoning for \(b\), we obtain two disjoint decompositions:
\[
	F_2 = S(a) \,\cup\, a S(a^{-1}),
	\qquad\text{and}\qquad
	F_2 = S(b) \,\cup\, b S(b^{-1}).
\]

Each of these equalities expresses \(F_2\) as a finite union of disjoint pieces
that are moved by a single group element (namely \(a\) or \(b\)) to reassemble the
whole group.  Taken together, they give two different reassemblies of \(F_2\)
from finitely many translated pieces.

\begin{corollary}\label{cor:F2_paradoxical-clean}
	The free group \(F_2\), acting on itself by left translation, admits a paradoxical
	decomposition.  In particular, \(F_2\) is \(F_2\)-paradoxical.
\end{corollary}


\begin{theorem}\label{thm:free_subgroup_SO3-updated}
	There exist rotations $A,B\in \mathrm{SO}(3)$ such that the subgroup
	\(
	H = \langle A,B\rangle
	\)
	is isomorphic to $F_{2}$.\end{theorem}
\begin{proof}[Sketch]
	Choose two orthogonal axes and a rotation angle
	\(\theta = \arccos(1/3)\).
	Let $A$ and $B$ be rotations by $\theta$ about these axes.
	A classical coordinate argument shows that no nontrivial reduced word in $A^{\pm1}, B^{\pm1}$
	fixes the point $(1,0,0)$, hence no reduced word represents the identity rotation.
	Thus $\langle A,B\rangle$ is a free subgroup on two generators.
\end{proof}

We henceforth identify $F_2$ with its copy
\(H = \langle A,B\rangle \subset \mathrm{SO}(3)\).

\subsection{A paradoxical decomposition of the sphere}
Let $S^2$ be the unit sphere.
The action of $H$ on $S^2$ is not free everywhere: rotations have two fixed points each.

\begin{lemma}\label{lem:fixed_points_countable_polished}
	Let $D\subset S^{2}$ be the set of points fixed by some nontrivial element of $H$. Then $D$ is countable.
\end{lemma}


\begin{proof}
	Every nontrivial rotation fixes exactly two antipodal points, and $H$ is countable.
	Thus $D$ is a countable union of two-point sets.
\end{proof}

Let \(S^2_0 := S^2 \setminus D\).
The action of \(H\) on \(S^2_0\) is free.
By the axiom of choice, we may select a set \(M \subset S^2_0\) containing exactly one
representative from each \(H\)-orbit.
The map
\[
	H\times M \to S^2_0,\qquad (h,m)\mapsto h\cdot m
\]
is bijective.
Let
\[
	H = \bigcup_{i=1}^{k} E_i
\]
be the finite partition of $H$. For each \(i\) define
\[
	E_i \cdot M := \{ h\cdot m : h\in E_i,\ m\in M\}.
\]

\noindent Then sets $E_i\cdot M$ form a partition of $S^2_0$.

\begin{proposition}\label{prop:S2_paradoxical_noD-updated}
	$S^2_0$ admits a paradoxical decomposition under the action of $H$.
\end{proposition}

\begin{proof}
	Apply the paradoxical decomposition of $H$ and transfer each piece $E_i$
	to $E_i\cdot M$.
	The relations are preserved because the action of $H$ on $S^2_0$ is free.
\end{proof}

\begin{proposition}\label{prop:S2_paradoxical_noD-updated}
	\(S^2_0\) admits a paradoxical decomposition under the action of \(H\).
\end{proposition}

\begin{proof}
	Transfer the paradoxical decomposition of \(H\) to \(S^2_0\) using the bijection
	\(H\times M\to S^2_0\).
	We obtain disjoint families
	\[
		A_i\cdot M,\qquad B_j\cdot M
	\]
	whose translates satisfy
	\[
		S^2_0=\bigcup_{i=1}^m g_i(A_i\cdot M)
		\qquad\text{and}\qquad
		S^2_0=\bigcup_{j=1}^n h_j(B_j\cdot M),
	\]
	which is a paradoxical decomposition of \(S^2_0\).
\end{proof}

To finish, we restore the omitted pole points; being countable, they do not affect
equidecomposability.

\begin{lemma}\label{lem:S2_equidecomp_fullsphere-updated}
	$S^2$ and $S^2_0$ are $H$-equidecomposable.
\end{lemma}

\begin{proof}
	$D$ is countable, so by Remark~\ref{rem:countable_sets harmless},
	removing $D$ does not affect equidecomposability.
	Thus
	\[
		S^2 \sim_H S^2 \setminus D = S^2_0.
	\]
\end{proof}

\begin{corollary}\label{cor:S2_paradoxical-updated}
	The unit sphere $S^2$ admits a paradoxical decomposition under rotations.
\end{corollary}

\subsection{From the sphere to the ball}

We now extend the paradoxical decomposition of the sphere to the closed unit ball
\(
B = \{x \in \mathbb{R}^3 : \|x\|\le 1\}.
\)

\begin{lemma}
	If \(S^2 = \bigcup A_i\) is paradoxical under rotations, then the sets
	\[
		\widetilde{A}_i := \{ r x : x\in A_i,\ 0<r\le 1 \}
	\]
	form a paradoxical decomposition of \(B\setminus\{0\}\).
\end{lemma}

\begin{proof}
	Every point of \(B\setminus\{0\}\) has a unique radial representation
	\(x = r u\) with \(u\in S^2\) and \(0<r\le1\).
	Rotations preserve radial lines:
	\(
	R(rx) = r\,R(x).
	\)
	Hence any paradoxical decomposition of \(S^2\) lifts along radial segments
	to a paradoxical decomposition of \(B\setminus\{0\}\).
\end{proof}

Finally we reincorporate the center point.

\begin{lemma}\label{lem:ball_center_equidecomp-updated}
	\(B\) and \(B\setminus\{0\}\) are equidecomposable.
\end{lemma}

\begin{proof}
	The removed set \(\{0\}\) is countable, so by
	Remark~\ref{rem:countable_sets harmless} we have
	\(
	B \sim B\setminus\{0\}.
	\)
\end{proof}

\begin{theorem}[Banach--Tarski Paradox]\label{thm:bt}
	The closed unit ball \(B\subset\mathbb{R}^3\) admits a paradoxical
	decomposition under rotations.
	In particular, \(B\) is equidecomposable with two disjoint copies of itself.
\end{theorem}

\begin{proof}
	By Corollary~\ref{cor:S2_paradoxical-updated}, the sphere \(S^2\) is paradoxical.
	Radial extension gives a paradoxical decomposition of \(B\setminus\{0\}\).
	By Lemma~\ref{lem:ball_center_equidecomp-updated},
	\(B \sim B\setminus\{0\}\).
	Since paradoxicality is preserved under equidecomposability
	(Lemma~\ref{lem:equidecomp_preserves_paradoxical}),
	the ball \(B\) itself is paradoxical.
\end{proof}


\newpage

\bibliographystyle{plainurl}
\bibliography{references}

\end{document}
